%============================ MAIN DOCU =================================
% define document class
\PassOptionsToPackage{table}{xcolor}
\documentclass[
   fontsize=10.5pt,
   invert-title=true,
   titlepage=false,
   titleimage-ratio=13,
   class=article,
   twocolumn,
   parskip=half-
]{bfhpub}				% KOMA-script report
%---------------------------------------------------------------------------
% Documents paths
%---------------------------------------------------------------------------
\makeatletter
\def\input@path{{content/}}
%or: \def\input@path{{/path/to/folder/}{/path/to/other/folder/}}
\makeatother
%-----------------  Set default font family  -------------------------------
\renewcommand{\familydefault}{\sfdefault}
%-----------------  Base packages     --------------------------------------
% Include Packages
\usepackage[english,ngerman,main=english]{babel}
\usepackage[babel, german=quotes]{csquotes}
%---------------------------------------------------------------------------
\usepackage{blindtext}  %% For placeholder text only
%---------------------------------------------------------------------------
\usepackage{geometry}
\geometry{
   a4paper, % Papierformat
   top=3cm, bottom=4cm, % Rand oben/unten
   outer=2.5cm, inner=3cm % aussen/inne
}
%---------------------------------------------------------------------------
\usepackage[
  backend=biber,
  style=numeric
]{biblatex}
\addbibresource{sample.bib} %% Name of bibliography file
%---------------------------------------------------------------------------
% Hyperref Package (Create links in a pdf)
%---------------------------------------------------------------------------
\usepackage[
   hidelinks,
   colorlinks,
   linkcolor=.,
   filecolor=.BFH-MediumGreen,
   urlcolor=BFH-MediumBlue,
   citecolor=.,
   plainpages=false,
   pdfpagelabels,
   pdfusetitle,
   hypertexnames = {true},	% no failures "same page(i)"
]{hyperref}
%---------------------------------------------------------------------------

\newcommand*{\code}[1]{\enquote{\texttt{#1}}}  %% The "code" macro definition
% chktex-file 41
\ohead*{\csname@author\endcsname{}}
\chead*{}
\ihead*{\csname@title\endcsname{}\title}

\setcounter{secnumdepth}{0}

\LoadBFHModule{boxes, terminal}

\begin{document}
%----------------  BFH tile page   -----------------------------------------
\title{Federated Learning for Wearable Stress-Prediction}
\subtitle{Project 2 Proposal}
\author{Cedric Aebi}
% \subject{BTF3232}
%\publishers{publishers}
% \institution{Bern University of Applied Sciences}
% \department{Technik und Informatik}
% \institute{Mikro- und Medizintechnik}
\version{1.0.0}
%The starred variant will automatically scale and clip the image. the non-starred one will allow you to set the size yourself


\makeatletter
\twocolumn
[
  \begin{@twocolumnfalse}
    \maketitle
    % \section{Abstract}
    \begin{itemize}
      \item \textbf{Institute} \textemdash{} Institute for Data Applications and
            Security (IDAS)
      \item \textbf{Degree Program} \textemdash{} Master of Science in
            Engineering with Specialization in Data Science
      \item \textbf{Advisor} \textemdash{} Dr.\ Souhir Ben Souissi
      \item \textbf{Date} \textemdash{} August 04, 2024
    \end{itemize}
  \end{@twocolumnfalse}
  \vskip1.5em]
\makeatother


%\section*{Abstract}
%\textbf{\blindtext[1]}


%\begin{multicols}{2}
\section{Introduction}
In today's fast-paced modern world, stress has become an unavoidable aspect of
daily life for many individuals. Whether it originates from the demands of work,
relationships, or personal challenges, its influence can deeply impact physical,
emotional, and cognitive well-being. Prolonged stress can lead to mental and
health problems such as depression or anxiety. Therefore, it is important to
monitor stress levels to minimize the negative effects of long-term stress. With
advancements in wearable technologies, smart gadgets like the Apple Watch can
help detect stress using their onboard sensors. Machine Learning can then be
applied to the collected data from these smart devices to predict stress levels.
However, for these machine learning algorithms to be effective, a significant
amount of data must be gathered from individuals. Not everyone is willing to
freely share their private and sensitive data. Therefore, in this study,
Federated Learning, which preserves privacy in machine learning, will be
analyzed and compared to traditional machine learning methods for predicting
stress using an already collected dataset on stress prediction.

\section{Related Work}
Talha Iqbal et al.~\cite{iqbal_stress_2022} used smartwatches (namely the
Empatica E4 watch) and created the dataset named `Stress-Predict Dataset' with
the help of 35 healthy volunteers. The data was collected using the smartwatches'
photoplethysmogram (PPG) sensor. The raw PPG signal was then filtered to
obtain a clean blood volume pulse (BVP) signal, which was then passed to an
estimation algorithm for heart rate (HR) and inter-beat intervals (IBI)
estimations. The volunteers underwent a series of tasks (i.e., Stroop color
test, Trier Social Stress Test, and Hyperventilation Provocation Test) to collect
data for the stress-predict dataset. In the same field of stress prediction
Fauzi et al.~\cite{fauzi_comparative_2022} compared three different learning
methods, \textit{Individual Learning}, \textit{Centralized Learning} and
\textit{Federated Learning}, on the WESAD (Wearable Stress and Affect Detection)
dataset~\cite{schmidt_introducing_2018}. They used a linear regression model in
all three learning methods and did the evaluation on the clients itself. Their
results showed, that Federated Learning showed lesser performance on different
metrics compared to individual or centralized learning.

\section{Research Question}
In this work (project 2), the author wants to take the same approach as Fauzi et
al.~\cite{fauzi_comparative_2022} and apply it to the stress predict dataset
from Talha Iqbal et al.~\cite{iqbal_stress_2022}. As a further step, instead of
linear regression models, more complex models such as DNN (Deep Neural Network)
can be adapted and applied to analyze the dataset. This leads to the general
research question:
\begin{displayquote}
  `How does Federated Learning perform compared to centralized learning and
  individual learning on a wearable stress-predict dataset?'
\end{displayquote}

\section{Materials and Methods}
\subsection{Dataset}
The stress-predict dataset\footnote{\url{https://github.com/italha-d/Stress-Predict-Dataset}} from Talha Iqbal et al.~\cite{iqbal_stress_2022} will
be used in this work. The dataset contains data from 35 healthy volunteers,
which performed three different stress-inducing tasks (i.e., Stroop colour word
test, interview session and hyperventilation session) with a baseline/relaxation
period between each task, for 60 minutes. Blood volume pulse (BVP),
Inter-beat-intervals, and heart rate were continuously recorded using Empatica
watches, while the respiratory rate was estimated using a PPG-based respiratory
rate estimation algorithm. The data is labeled for stress and non-stress.

\subsection{Comparing Approaches}
For this study three different machine learning approaches will be compared:
\begin{itemize}
  \item Centralized Learning
  \item Federated Learning
  \item Individual Learning
\end{itemize}

In the \textit{Individual Learning} scenario, the dataset gets split into
multiple different parts to simulate different clients. For each part, one
separate model will be trained on the available data in this part. In the
\textit{Federated Learning} scenario, the simulated clients will participate in
collective training. For this, a publicly available FL framework like TensorFlow
Federated\footnote{\url{https://www.tensorflow.org/federated}}
or Flower\footnote{\url{https://flower.ai}} will be used. In the \textit{Centralized Learning}, one model gets trained
on all the available data in the dataset. This scenario serves as a direct
comparison to FL, to see if FL loses any performance, while maintaining its
strengths like data privacy.

Different models (linear regression, simple neural networks etc.) will be used
and compared in the different learning scenarios.

\subsection{Evaluation}
The datasets will be split into parts to simulate individual clients, and then
each client will perform an 80\%-20\% train-test split on its available data. As
evaluation metrics, the F1 score, precision, recall, and accuracy will be
used. For the federated learning approach, both client-side evaluation (the
performance will be measured on the client-side using its test data) and
server side evaluation (the
performance will be measured on the server-side using its global held test data) will be done.

\section{Expected Results}
Theoretically, since FL has the ability to access all the available data from all
the clients but without sharing the raw data, it should show similar
performance (F1 score, precision, recall and accuracy) compared to the
centrally trained approach. The locally trained models should perform
significantly worse compared to both the FL and centrally trained approaches
since the local devices have only a fraction of the available data, leading the models
to underfit. However, Fauzi et al.~\cite{fauzi_comparative_2022}
showed in their paper that FL underperformed on the WESAD dataset. So it's hard
to say which results will come out. A well-fitting model has to be chosen.

\section{Milestones}
\begin{enumerate}
  \item April 12, 2024 \textemdash{} Hand-in final proposal
  \item May 05, 2024 \textemdash{} First implementation done with simple
        linear regression model
  \item June 09, 2024 \textemdash{} Final implementations with different models
        and different evaluation techniques
  \item July 12, 2024 \textemdash{} Project 2 Presentation
  \item August 04, 2024 \textemdash{} Hand-in final paper
\end{enumerate}

%\end{multicols}

%------------ BIBLIOGRAPHY ---------------

\onecolumn
\printbibliography{}


\end{document}
